\section{问题一：交会定位区域的直径算法与覆盖性判定}\label{sec:q1}

\subsection{从示向度读数到线性约束}

设在检测点 $s_k$ 处测得某干扰源的示向度为 $\theta_k$，误差界为 $\varepsilon$。由\cref{asu:bounded}，真源 $g$ 相对 $s_k$ 的真实方位角落在 $[\theta_k-\varepsilon,\theta_k+\varepsilon]$ 内，因此 $g$ 位于以 $s_k$ 为顶点、张角 $2\varepsilon$ 的\textbf{前向楔形}
\begin{equation}\label{eq:wedge}
  W_k=\Bigl\{x\in\R^2 \;\Bigm|\; \bigl[\uvec{\theta_k-\varepsilon},\,x-s_k\bigr]\ge 0,\quad
                                  \bigl[\uvec{\theta_k+\varepsilon},\,x-s_k\bigr]\le 0\Bigr\}
\end{equation}
之中，其中 $[\bm a,\bm b]=a_xb_y-a_yb_x$。定位区域即
\begin{equation}\label{eq:P}
  P=\bigcap_{k=1}^{K} W_k .
\end{equation}

\begin{lemma}[两条不等式已自动选出前向分支]\label{lem:forward}
当 $2\varepsilon<180^\circ$ 时，\cref{eq:wedge}恰好刻画以 $\uvec{\theta_k}$ 为轴、半角为 $\varepsilon$ 的单侧锥，而不包含背向方向。
\end{lemma}
\begin{proof}
令 $x-s_k=r\,\uvec{\varphi}$（$r>0$）。两条约束分别化为 $\sin(\varphi-\theta_k+\varepsilon)\ge0$ 与 $\sin(\varphi-\theta_k-\varepsilon)\le0$，即
$\varphi-\theta_k+\varepsilon\in[0,\pi]$ 且 $\varphi-\theta_k-\varepsilon\in[-\pi,0]\pmod{2\pi}$，两者之交为 $\varphi-\theta_k\in[-\varepsilon,\varepsilon]$。特别地取 $\varphi=\theta_k+180^\circ$ 时第一式给出 $-\sin\varepsilon<0$，故背向射线被排除。
\end{proof}

\cref{lem:forward}并非可有可无的细节：若把示向度处理成一条“带宽为 $2\varepsilon$ 的双向直线带”，则背向的一整块区域会被错误地保留，在多站交会时可能产生虚假的交集。

把\cref{eq:wedge}的两条约束写成行向量形式，$K$ 个检测点共给出 $m=2K$ 条线性不等式，于是
\begin{equation}\label{eq:poly}
  P=\{x\in\R^2 \;:\; \bm A x\le \bm b\},\qquad \bm A\in\R^{m\times2},\ \bm b\in\R^{m}.
\end{equation}
$P$ 是闭凸集，但\textbf{不一定有界}，也可能为空或退化。

\subsection{直径的计算：算法与正确性}

\begin{theorem}[直径在顶点上取到]\label{thm:diam}
设 $P$ 为有界闭凸多边形，顶点集为 $V(P)=\{v_1,\dots,v_k\}$，则
$\diam(P)=\max_{i<j}\|v_i-v_j\|$。
\end{theorem}
\begin{proof}
$\|\cdot\|$ 是凸函数，故对固定的 $y$，$x\mapsto\|x-y\|$ 在紧凸集 $P$ 上的最大值必在极点处取到；对 $y$ 再用一次同样的论证即可。由 Minkowski 定理，有界闭凸多边形的极点恰为其顶点。
\end{proof}

由\cref{thm:diam}，求直径只需枚举顶点对，而不必在边界上做连续优化。据此给出\cref{alg:diam}。它的关键不在于求直径这一步本身，而在于\textbf{先把四类情形分清楚，再计算}：直接对约束做顶点枚举而不判空、不判无界，会在读数互相矛盾或测站张角不足时输出毫无意义的数字。

\begin{algorithm}[!htbp]
\caption{有界误差交会定位区域的直径算法}\label{alg:diam}
\begin{algorithmic}[1]
\Require 检测点 $s_1,\dots,s_K$，示向度 $\theta_1,\dots,\theta_K$，误差界 $\varepsilon$
\Ensure 定位区域 $P$ 的状态（空／无界／有界）、顶点表与直径 $D$
\State 由\cref{eq:wedge}组装 $\bm A\in\R^{2K\times2}$，$\bm b\in\R^{2K}$
\State 解可行性线性规划 $\min\{0 \;:\; \bm Ax\le\bm b\}$
\If{不可行} \Return \textbf{空集}（读数与误差界相互矛盾，应报错而不是输出伪定位点）
\EndIf
\For{$\bm c\in\{(1,0),(-1,0),(0,1),(0,-1)\}$}
  \State 解 $\min\{\bm c^{\!\top}x \;:\; \bm Ax\le\bm b\}$
  \If{无下界} \Return \textbf{无界}，$D=+\infty$
  \EndIf
\EndFor
\State $V\gets\varnothing$
\For{每一对约束 $(i,j)$，$i<j$}
  \If{$\det(\bm a_i,\bm a_j)\ne0$}
    \State 解 $\bm a_i^{\!\top}x=b_i,\ \bm a_j^{\!\top}x=b_j$ 得交点 $x$
    \If{$\bm Ax\le\bm b+\tau\bm 1$ 且 $x\notin V$（按容差去重）} $V\gets V\cup\{x\}$
    \EndIf
  \EndIf
\EndFor
\State 以 $V$ 的重心为极点，按极角对 $V$ 排序，得到凸多边形顶点序列
\State $D\gets\max_{v_i,v_j\in V}\|v_i-v_j\|$ \Comment{线段与单点自然退化为 $D=\|v_1-v_2\|$ 与 $D=0$}
\State \Return \textbf{有界}，$V$，$D$
\end{algorithmic}
\end{algorithm}

\paragraph{五次线性规划各自的用途.} 需要说明的是，算法中的线性规划\emph{只用来判别 $P$ 的类型，不用来求顶点}。第 2 行的可行性规划取零目标，只问“约束组是否相容”；第 5 至 7 行的四次规划分别沿 $\pm x$、$\pm y$ 方向求极小，只要有一次报告目标无下界，$P$ 就在该方向上延伸至无穷，属于无界情形。四个坐标方向足以判定：二维多面体无界当且仅当其回收锥非平凡，而非平凡锥必与某个坐标半轴张成正夹角。判完类型之后，有界情形的顶点由第 9 至 14 行的“边界线两两求交、再逐个检验是否满足全部约束”得到，与线性规划无关。

\paragraph{复杂度.} 上述实现的开销为：$5$ 次二维线性规划（规模均为 $m$ 条约束、$2$ 个变量，可视为常数次）；顶点枚举 $O(m^2)$ 对交点，每个交点需 $O(m)$ 次可行性检验，合计 $O(m^3)$；最后直径枚举 $O(k^2)$。本题中单个源的检测点很少（$K\le4$，即 $m\le8$），这一实现完全够用，且代码短、边界情形容易查错。若检测点数很大，可换成半平面交（$O(m\log m)$）加旋转卡壳（$O(k)$）把总复杂度降到 $O(m\log m)$——这是一个\emph{替代实现}，本文提交的代码用的仍是前一种。

\paragraph{四类退化的处理.} $P$ 为空意味着读数与误差界不相容（例如误差界设小了），此时必须报错并排查角度约定、取整或坐标系，\emph{不能}静默地取一个近似交点；$P$ 无界出现在各楔形的方向锥未能互相封闭时（典型情形是只有一个检测点，或多个检测点近似共线且同侧），此时直径为 $+\infty$，用人为的大方框截断再报一个有限直径是错误的；$P$ 退化为线段或单点时，\cref{alg:diam}给出的 $D$ 自然正确，无需特判。

\subsection{以定位区域直径为直径的圆能否覆盖定位区域}

\begin{theorem}[覆盖性的充要条件与紧界]\label{thm:jung}
设 $P\subset\R^2$ 为有界闭凸集，$D=\diam(P)$，$\mecr=\mecr(P)$ 为最小包围圆半径。则
\begin{equation}\label{eq:jung}
  \frac{D}{2}\ \le\ \mecr\ \le\ \frac{D}{\sqrt3},
\end{equation}
右端为 Jung 不等式，两端均可取等。于是
\begin{itemize}
  \item 存在以 $D$ 为直径的圆覆盖 $P$，\textbf{当且仅当} $\mecr=D/2$；
  \item 一般情况下只能保证：以 $\dfrac{2D}{\sqrt3}\approx1.1547\,D$ 为直径的圆必能覆盖 $P$，且该系数不可改进。
\end{itemize}
\end{theorem}
\begin{proof}
左端：取 $p,q\in P$ 使 $\|p-q\|=D$，任何包含 $P$ 的圆都包含 $p,q$，故其半径不小于 $D/2$。右端即平面情形的 Jung 定理。取等：$P$ 为线段时 $\mecr=D/2$；$P$ 为边长 $D$ 的正三角形时 $\mecr=D/\sqrt3$。第一条由 $\mecr\ge D/2$ 立得。
\end{proof}

\paragraph{结论：一般不能覆盖.} 由\cref{thm:jung}，只要 $\mecr>D/2$，以 $D$ 为直径的圆无论圆心放在哪里都无法覆盖 $P$。这不是抽象反例，而是可以由真实的检测点与示向度产生的：

\begin{example}[三站交会给出 Jung 界上的正三角形]\label{exa:tri}
取三个检测点与读数
\[
  \begin{aligned}
  s_1&=(-1100,\,0), & \theta_1&=1^\circ;\\
  s_2&=(590,\,-952.6279), & \theta_2&=121^\circ;\\
  s_3&=(570,\,987.2690), & \theta_3&=241^\circ,
  \end{aligned}
\]
误差界 $\varepsilon=1^\circ$。按\cref{alg:diam}计算，$P$ 为有界三角形，顶点为
$(0,0)$、$(40,0)$、$(20,\,34.641016)$，即边长 $40\mmet$ 的正三角形。此时
\[
  D=40.000000\mmet,\qquad \mecr=23.094011\mmet,\qquad \frac{D}{2}=20\mmet,
\]
$\mecr/D=0.577350=1/\sqrt3$，恰好取到 Jung 不等式的等号。以 $D$ 为直径的圆（半径 $20\mmet$）\textbf{不能}覆盖该三角形，缺口为 $3.094011\mmet$，见\cref{fig:q1-counter}。
\end{example}

\begin{figure}[!htbp]
  \centering

  %% 画法建议：左子图画\cref{exa:tri}的三个检测点（远处，可用折断坐标轴或局部放大框）及三条示向度中心线与各自 $\pm1^\circ$ 的虚线，
  %% 交出的正三角形用浅色填充；右子图放大该三角形，叠加两个圆：半径 $D/2=20\mmet$ 的圆（红色虚线，可见三角形三个角伸出圆外）
  %% 与最小包围圆 $\mecr=23.094\mmet$（蓝色实线，恰好过三个顶点）。标注 $D=40\mmet$ 的一条边与 $\mecr/D=1/\sqrt3$。
  \fbox{\parbox[c][2.2cm][c]{0.86\textwidth}{\centering\small
  【待插图】（绘图要求见本段 \TeX{} 源码注释）}}
  \caption{以定位区域直径为直径的圆一般不能覆盖定位区域}\label{fig:q1-counter}
\end{figure}

\paragraph{什么时候可以覆盖：两站交会的情形.} 题目附图所示的\emph{两站}交会四边形，在两站到源的距离远大于区域尺度时趋于平行四边形；对这一类区域，同直径圆总是够用的。

\begin{lemma}[平行四边形的最小包围圆]\label{lem:para}
设 $P$ 为平行四边形，顶点为 $p,\ p+\bm u,\ p+\bm u+\bm v,\ p+\bm v$，则 $\mecr(P)=\diam(P)/2$，即以其直径为直径的圆一定能覆盖 $P$。
\end{lemma}
\begin{proof}
不妨设 $\bm u^{\!\top}\bm v\ge0$（否则把 $\bm v$ 换成 $-\bm v$ 并平移顶点标号，图形不变）。两条对角线的长度为 $\|\bm u\pm\bm v\|$，而
\[
  \|\bm u+\bm v\|^2-\|\bm u-\bm v\|^2=4\,\bm u^{\!\top}\bm v\ \ge 0,
\]
故 $\bm u+\bm v$ 是较长的对角线；又 $\|\bm u+\bm v\|^2=\|\bm u\|^2+\|\bm v\|^2+2\bm u^{\!\top}\bm v\ge\max(\|\bm u\|^2,\|\bm v\|^2)$，说明它也不短于任何一条边，于是 $\diam(P)=\|\bm u+\bm v\|$。

取圆心 $c=p+(\bm u+\bm v)/2$、半径 $\|\bm u+\bm v\|/2$。该圆过 $p$ 与 $p+\bm u+\bm v$；对另外两个顶点，
\[
  \bigl\|(p+\bm u)-c\bigr\|=\bigl\|(p+\bm v)-c\bigr\|=\frac{\|\bm u-\bm v\|}{2}\ \le\ \frac{\|\bm u+\bm v\|}{2},
\]
故四个顶点全在圆内或圆上，凸性给出 $P$ 被该圆覆盖，即 $\mecr(P)\le\diam(P)/2$。结合\cref{thm:jung}的左端不等式即得等号。
\end{proof}

\cref{lem:para}的论证只用到向量内积的符号，不需要对交会角作任何限制——这一点很重要，因为\cref{tab:q1-example}的第一个算例交会角为 $97.13^\circ$，已经是钝角。\cref{tab:q1-example}中 $\mecr$ 与 $D/2$ 在四位小数下完全相同，与引理一致。

\begin{table}[!htbp]
  \caption{两站交会的数值算例（真源 $g=(600,400)$，$\varepsilon=1^\circ$，读数含 $+0.6^\circ$ 与 $-0.4^\circ$ 的偏差）}
  \label{tab:q1-example}\centering\small
  \begin{tabular}{cccccccc}
    \toprule[1.5pt]
    $s_1$ & $s_2$ & 交会角 $\psi$ & 顶点数 & $D/\mmet$ & $\mecr/\mmet$ & $D/2$ & 可否一次认证清除 \\
    \midrule[1pt]
    $(0,0)$ & $(800,0)$  & $97.13^\circ$ & $4$ & $32.0756$ & $16.0378$ & $16.0378$ & 可（$\mecr\le20$） \\
    $(0,0)$ & $(1400,0)$ & $60.26^\circ$ & $4$ & $55.7265$ & $27.8632$ & $27.8632$ & 否，需光学兜底 \\
    \bottomrule[1.5pt]
  \end{tabular}
\end{table}

\noindent 两例中真源均满足全部线性约束，即 $g\in P$，符合\cref{asu:bounded}下的包含性要求。

\subsection{本问结论对后两问的决定性影响}

\cref{thm:jung}直接给出了后两问中\textbf{“何时可以一次性把源清掉”}的正确判据。由于清除成功的条件是清除点与源的距离不超过 $r_c=20\mmet$，而算法只知道 $g\in P$，故：

\begin{corollary}[认证清除判据]\label{cor:certify}
在 $c_*(P)$ 处执行一次 \texttt{/clear} 必定成功，当且仅当 $\mecr(P)\le r_c=20\mmet$。若仅已知直径，则
\begin{equation}\label{eq:diamsuff}
  D\le \sqrt3\,r_c=20\sqrt3\approx34.641\mmet
\end{equation}
是 $\mecr(P)\le r_c$ 的充分条件，且由\cref{exa:tri}该系数不可放宽到 $2r_c=40\mmet$。
\end{corollary}

因此，工程上常见的“定位区域直径小于 $2\times20=40\mmet$ 就能一次清除”的说法是\textbf{错误}的：\cref{exa:tri}中 $D=40\mmet$ 却存在距圆心 $23.09\mmet$ 的可能位置。本文后续算法一律以 $\mecr(P)\le20\mmet-10^{-6}$ 作为认证条件（保留浮点余量），而不使用直径判据；只有在需要一个只依赖直径的快速筛选时，才使用\cref{eq:diamsuff}。
